% Appendix F (writer: W5 of the restructure; plan: research/paper-review/restructure/PLAN.md). Proofs for sec:teacher.
% Sources: base app-teacher.tex §C.3-C.7 (proof of thm:teacher:main, numbers, full table, finite samples, ex:teacher:version1,
% prop:teacher:concentration); base app-limits.tex §D.1 (proof of prop:limits:plain); research/tracks/split/teacher/notes-final.md
% Prop 2.1, Remark 2.6, Lemma 3.0, Props 3.1, 3.2(a), 3.6, 3.7, Lemma 5.1, Props 5.2, 5.3; split/verify-teacher.md §2-§9
% (g >= 0; the corrected R_max; continuity of Phi_2^sharp; strict monotonicity of psi_2 and of psi_2(p,1)-H(p)).
% The proofs of the four lemmas and of the sharp floor (base §C.1-C.2) are in app-parity.tex (W4).
% prop:teacher:frozen, prop:limits:plain and tab:teacher:rates are stated here, with summaries in teacher.tex (page budget; log-W5);
% prop:teacher:sharp and its display moved here in the fix round (fix-W5 RW5-04), with a one-paragraph pointer in sec:teacher:rates.
% Integration of 2026-10-14 (research/paper-review/imitation-edit/log-integrate.md): app:teacher:weightdecay (prop:teacher:nofixedprior,
% ex:teacher:weightdecay) before app:teacher:rates, text of research/tracks/imitation/notes-final.md §3.3 (= drafts/wd-appendix.tex, with F12).
% Fix round of 2026-10-14 (research/paper-review/imitation-edit/log-fix.md): imitation-edit review-math L3; review-front m8 (omega_d for the decoupled coefficient), m9 (app:teacher:weightdecay).
\section{Proofs for \texorpdfstring{\cref{sec:teacher}}{Section \ref*{sec:teacher}} (the teacher's mistakes)}\label{app:teacher}

Throughout, the setting is \cref{def:teacher:instance}, and (F1)--(F5) are the facts of \cref{app:setup:facts}. For $z\in\{0,1\}^m$ and $w\in\{0,1\}^{d-m}$, $(z,w)$ is the input with prefix $z$, and $w$ is uniform in every $\Ex_w$. We write $\tau_B:=2^{1-B}$ ($0$ if $B=\infty$), so that $G_j=1-\log(1+2^{-j}+2p+\tau_B)$; $T_g:=\delta_S+\Delta_\ell+\Lambda_3+1+2\log m+2\log s_0+g+2\log g+\kappa$, the threshold that defines $j_g$ in \cref{thm:teacher:main}(ii)(e); $\mathrm{loss}_\pi(x):=-\log\nu_\pi(y^*(x)\mid x)\in[0,\infty]$; and $N_n(x)$ for the number of queries at $x$ among the first $n$.

\subsection{Noise and errors: frozen noise, plain Bayes and schedules}\label{app:teacher:schedules}

\begin{proposition}[Frozen noise, counted once or per query]\status{proved}\src{split/teacher Prop 3.7; split/verify-teacher \S7}\label{prop:teacher:frozen}
Let the inputs be i.i.d.\ from a law $\mu$ on $X$, let $\nu$ be a sampler with a program $\pi_0$ of length $k$, in a class $\mathcal C$ as in \cref{def:split:setting}, whose kernel equals $\nu$ on $\supp\mu$, and let the table $F$ have independent values $F(x)\sim\nu(\cdot\mid x)$, drawn once.
(a) \emph{FR.} For fixed $\lambda>1$, at every $n$, repeats or not, every $\lambda$-MAP has length at most $k+(k+\delta^*)/(\lambda-1)$, where $\delta^*$ is the Ville deficiency of $\xi_{\mathcal C}$ along the distinct inputs (\cref{lem:split:ville}), so $\Prob[\delta^*\ge j]\le2^{-j}$. A program that agrees with $\nu$ off a finite set $A\subseteq\supp\mu$ fits $D_n$ better than $\pi_0$ by at most $\sum_{x\in A}\log(1/\nu(F(x)\mid x))$ bits, however often the inputs of $A$ are queried.
(b) \emph{FQ.} Let $A\subseteq\supp\mu$ be finite with $\nu(F(x)\mid x)<1$ for some $x\in A$, and let $\pi_A\in\mathcal C$ agree with $\nu$ off $A$ and give $F(x)$ probability $1$ on $A$. The FQ fits of $\pi_0$ and $\pi_A$ differ by $\sum_{x\in A}N_n(x)\log(1/\nu(F(x)\mid x))$. So for every schedule $\lambda_n=o(n)$, almost surely each program whose kernel equals $\nu$ on $\supp\mu$ is a $\lambda_n$-MAP for only finitely many $n$.
(c) If $\mu$ is uniform on a finite set $X_0$ and every input of $X_0$ has been queried $\bar n$ times, the FQ $\lambda$-MAPs are the $(\lambda/\bar n)$-MAPs of the table $F|_{X_0}$, so every $\lambda$-threshold of that table is multiplied by $\bar n$.
\end{proposition}

By (c) and \cref{cor:teacher:fulltable}, in FQ frozen noise needs $\lambda$ of order $n/|X_0|$, like errors with $m=d$.

\begin{proof}
(a) The next new input is fixed by the i.i.d.\ input stream, independent of $F$, and $F(x)$ at a new $x$ is independent of the earlier labels and inputs; almost surely all queried inputs lie in $\supp\mu$, so $\nu_{\pi_0}(D_n)=\nu(D_n)$ for all $n$. So \cref{thm:split:private}(i) applies with $\pi_\nu=\pi_0$, and \cref{lem:split:ville} gives the tail of $\delta^*$. If $\pi$ agrees with $\nu$ off $A$, then $\nu_\pi(D_n)/\nu_{\pi_0}(D_n)=\prod_{x\in A\cap\dom D_n}\nu_\pi(F(x)\mid x)/\nu(F(x)\mid x)\le\prod_{x\in A}1/\nu(F(x)\mid x)$.
(b) Off $A$ the per-query losses agree; on $A$, $\pi_A$ has loss $0$ and $\pi_0$ loss $\log(1/\nu(F(x)\mid x))$ at each query. As $X$ is countable, $\supp\mu$ is the set of atoms, and $N_n(x)/n\to\mu(x)>0$ almost surely; so the difference divided by $n$ tends to a positive limit. Every program $\pi'$ whose kernel equals $\nu$ on $\supp\mu$ has the FQ fit of $\pi_0$, so $J_{\lambda_n}(\pi')-J_{\lambda_n}(\pi_A)\ge-\lambda_n\ell(\pi_A)+\sum_{x\in A}N_n(x)\log(1/\nu(F(x)\mid x))\to+\infty$ when $\lambda_n=o(n)$.
(c) The FQ fit is $\bar n\,L(\pi;F|_{X_0})$; apply \cref{lem:teacher:reuse}.
\end{proof}

\begin{proposition}[Plain Bayes, MDL and SI learn the errors]\status{proved; computed}\src{steeper Props 2.9, 2.10; checks c3, c6(G), c7(G)}\label{prop:limits:plain}\label{app:limits:plain}
(a) Let $P$ be a shared hypothesis with $P\ge2^{-\ell(\pi)-\kappa}\nu_\pi$ for every $\pi\in\Hcl{\times}$, such as $\Mprod$ or $\MF$. For every sequence of distinct inputs and every $n$, $\sum_{i\le n}-\log P(y^*(x_i)\mid x_i,D_{i-1})\le\ell(\pi_b)+\kappa\le d+\lceil L_S\rceil+\kappa_b+\kappa$. So at most $d+\lceil L_S\rceil+\kappa_b+\kappa$ inputs in total ($d+L_S+O(\log d+\log m)$ if $p=2^{-r}$) get probability below $\frac12$ for the teacher's label (semimeasure convention, hence also normalised), and every other revealed error input gets the erroneous label with probability at least $\frac12$.
(b) Let $B<\infty$, $\delta\in(0,1)$, let $\lambda=\lambda_n\ge1$ be fixed or depend on $n$, and draw $n\le2^d$ inputs uniformly without replacement. With probability at least $1-\delta$, every $\lambda$-MAP $\hat\pi$ has $R_B(\hat\pi)\le\lambda\hat L/n+B\sqrt{(\hat L\ln2+\ln(2/\delta))/(2n)}$ with $\hat L:=d+\lceil L_S\rceil+\kappa_b$, hence $e(\hat\pi)\ge2^{-R_B(\hat\pi)/p}-2^{-B}$. In particular $e(\hat\pi)\ge1-2^{-B}-o(1)$ whenever $(\lambda_n/n)\hat L\to0$ and $\hat L/n\to0$, along any sequence of instances $(d_n)$ with fixed $m$, $p$ and error rule and $n\le2^{d_n}$. On a fixed instance the bound is informative only once $(\lambda_n/n)(d+L_S)$ is small compared with $p$; for slowly decaying schedules such as $\lambda_n=n^{0.9}$ at the example's sizes this never happens for $n\le2^d$.
(c) On the interleaved string with i.i.d.\ uniform inputs, Solomonoff's $M$ has expected total label loss at most $d+\lceil L_S\rceil+\kappa_b+\kappa$, and for fixed $\lambda$ the steeper SI predictor pays at most $\lambda\ell_U$ more (\cref{thm:absorb:absorption}(d)).
\end{proposition}

For instance, at the example's sizes with all inputs revealed and the machine constant set to $0$ (an illustration), a direct computation (code in the accompanying repository) shows that the exceptions in (a) are at most $11{,}853$ of the $1.238\cdot10^{27}$ error inputs, a fraction at most $9.6\cdot10^{-24}$. With $B=30$ and $\delta=0.01$ (machine constant $0$, an illustration), the same computation gives $e\ge0.649$ at $n=10^{13}$ for $\lambda=1$, and $e\ge0.734$ at $n=10^{15}$ for $\lambda_n=\sqrt n$.

\begin{proof}
(a) The memoriser is consistent with every $D_n$, so $P(D_n)\ge2^{-\ell(\pi_b)-\kappa}$, and the sequential loss of a shared hypothesis telescopes to $-\log P(D_n)+\log P(\emptyset)\le\ell(\pi_b)+\kappa$. Each term is nonnegative, a probability below $\frac12$ costs more than a bit, and normalised conditionals are at least the raw ones. For $p=2^{-r}$ the memoriser can name $s_0=2^{m-r}$ through $r$. For $\MF$, $P\ge2^{-K(\nu_\pi)-O(1)}\nu_\pi$ as product laws are machine laws, and $K(\nu_\pi)\le\ell(\pi)+\kappa$; this is \citet[Corollary~5.4]{claudeSFI2026fundamentals}.% P1: cor:learn:mistakes
\
(b) Apply \cref{lem:teacher:finite} with $\pi_{\rm ref}=\pi_b$, so $A=0$, $\epsilon_1=0$, $c=\lambda_n/n$ and $\hat L=\ell(\pi_b)$; it holds for each $n$, so $\lambda$ may depend on $n$. The error-label bound is the last step of the proof of \cref{thm:teacher:main}(iii). The bound on $R_B$ is at least $(\lambda_n/n)(d+L_S)$; for $\lambda_n=n^{0.9}$ at the example's sizes, $\lambda_n/n\ge2^{-10}=p$ for every $n\le2^{100}$.
(c) The law $\Prob$ of the interleaved sequence (fixed-length blocks) is computable from $a$, $S$, $d$, $m$ and $s_0$, so it has a program in $\Hcl{\mathrm{SI}}$ of length at most $d+\lceil L_S\rceil+\kappa_b+\kappa$, and $M$ dominates it; the chain-rule argument in the proof of \cref{thm:absorb:absorption}(d) bounds the expected label loss by that length. For $\lambda>1$ apply \cref{thm:absorb:absorption}(d), with $\kappa_0=0$ in $\Hcl{\mathrm{SI}}$.
\end{proof}

\begin{proof}[Proof of \cref{prop:teacher:schedules}]
Let $\hat R_n(\pi):=\sum_{x\in\{0,1\}^d}(N_n(x)/n)\,\mathrm{loss}_\pi(x)$. As $N_n(x)/n\to2^{-d}$ for each of the finitely many $x$, almost surely $\hat R_n(\pi)\to R(\pi)$ for every $\pi$ (if $R(\pi)=\infty$, $\hat R_n(\pi)=\infty$ once an input of infinite loss is queried).
(a) The memoriser has fit $0$, so a $\lambda_n$-MAP $\hat\pi$ has $\lambda_n\ell(\hat\pi)\le\lambda_n\ell(\pi_b)$; as $\lambda_n>0$ and by Kraft's inequality, finitely many candidates remain. For each with $R(\pi)>0$, $\frac1n[J_{\lambda_n}(\pi)-J_{\lambda_n}(\pi_b)]=\frac{\lambda_n}n(\ell(\pi)-\ell(\pi_b))+\hat R_n(\pi)\to R(\pi)>0$, so it is a $\lambda_n$-MAP for only finitely many $n$. Under uniform inputs $R(\pi)=0$ means $\nu_\pi(y^*(x)\mid x)=1$ at every $x$, so $e(\pi)=1$.
(b) The coin has $R=1$, so $\ell_{\min}$ exists; fix $\pi_0$ with $\ell(\pi_0)=\ell_{\min}$ and $R(\pi_0)<\infty$. Lengths are integers, so every $\pi$ with $\ell(\pi)>\ell_{\min}$ has $J_{\lambda_n}(\pi)\ge\lambda_n(\ell_{\min}+1)>J_{\lambda_n}(\pi_0)$ once $\lambda_n>n\hat R_n(\pi_0)$, which holds eventually. In $\Hcl{\times}$, $\ell_{\min}\le\ell(\pi_{\rm coin})\le\kappa$, while level-$1$ programs have length at least $d-4-\Lambda_1$ by (D1).
(c) Fix $n$ and $c:=\lambda_n/n\in[c_1,c_2]$; a $\lambda_n$-MAP minimises $\hat J_c:=J_{\lambda_n}/n$. \Cref{lem:teacher:finite} with $\pi_{\rm ref}=\pi_g$ (losses in $[0,\log(1/p)]$) holds for every $c$, because its two events do not involve $c$; for PQ data, Hoeffding's inequality for independent bounded variables gives the same events. Its excess has $\hat L(c)=(J_c(\pi_g)+\epsilon_1)/c\le\bar L$, as $J_g$ is increasing and $c\ge c_1$. So with probability at least $1-\delta$, $J^B_c(\hat\pi)\le J_g(c)+\bar\epsilon_n$, and \cref{thm:teacher:main}(ii) applies. $\Phi_2-J_g$ is a minimum of finitely many strictly decreasing affine functions, positive on $[c_1,c_2]\subset(0,c_{\rm hi})$, so it has a positive minimum there, while $\bar\epsilon_n\to0$, also with $\delta=n^{-2}$; then the failure probabilities are summable, and the Borel--Cantelli lemma gives the almost-sure form.
\end{proof}

\subsection{The kink bounds and the theorem}\label{app:teacher:main}

Write $\psi_c(R):=c\max\{0,B_S-2^mR\}+R=\max\{R,\ cB_S-(c2^m-1)R\}$. We prove (o), then \cref{prop:teacher:kinks}, then the rest.

\begin{proof}[Proof of \cref{thm:teacher:main}(o)]
If $\lev(\pi)=1$, (D1) and (D2) with $s(1)=4$ give $\ell(\pi)\ge d-4-\Lambda_1+\max\{0,B_S-2^mR_B(\pi)\}$; multiply by $c$ and add $R_B(\pi)$. If $\lev(\pi)=t\in[2,\infty)$, (D1), $\ell\ge0$ and \cref{lem:teacher:floor} give $J^B_c(\pi)\ge c(d-s(t)-\Lambda_1)^++G_{t-1}\ge\varphi(t):=c(d-4t-\Lambda_1)^++G_{t-1}$, as $s(t)\le4t$. Extend $\varphi$ to real $t\ge2$. If $t^*\ge2$: on $[2,t^*]$ the first term is linear and $G_{t-1}$ concave, so $\varphi\ge\min\{\varphi(2),\varphi(t^*)\}$; for $t>t^*$, $\varphi(t)\ge G_{t-1}\ge G_{t^*-1}$. If $t^*<2$: $\varphi(t)\ge G_1$. At level $\infty$, $J^B_c\ge R_B\ge G_\infty\ge G_j$ for every $j$. ($G_j$ is increasing and concave in $j$: with $v=2^{-j}$, $\mathrm dG/\mathrm dj=v/(1+v+2p+\tau_B)>0$ and $\mathrm d^2G/\mathrm dj^2=-\ln2\cdot v(1+2p+\tau_B)/(1+v+2p+\tau_B)^2<0$.)
\end{proof}

\begin{proof}[Proof of \cref{prop:teacher:kinks}]
(a) By (o), every $\pi\in\mathcal N$ has $J^B_c(\pi)\ge\Phi_2(c)$, and the reference value is $J_g(c)$; $J_g-\Phi_2$ is a maximum of strictly increasing affine functions, hence $\le0$ exactly on $(0,c_{\rm hi}]$ and $<0$ on $(0,c_{\rm hi})$. At $c=c_{\rm hi}>0$, $c_{\rm hi}\ell(\pi)+R_B(\pi)\ge\Phi_2(c_{\rm hi})\ge J_g(c_{\rm hi})$, the slanted side of the wedge; the flat side is $R_B(\pi)\ge G_{\lev(\pi)-1}\ge G_1$ (\cref{lem:teacher:floor}).
(b) (D1) with $s(1)=4$ gives $\ell(\pi)\ge d-4-\Lambda_1$, and (D2) gives $R_B(\pi)\ge B_S/2^m-2^{-m}(\ell(\pi)-d+4+\Lambda_1)=H(p)-\Lambda_2/2^m-2^{-m}(\ell(\pi)-d-\kappa_g)$. If $p<\frac14$, then $u(\pi_g)=1-p>\frac34$, so $\pi_g\in\mathcal M$, and the line at $\pi_g$ ($\ell(\pi_g)\le d+\kappa_g$, $R_B(\pi_g)=H(p)$) gives $\Lambda_2\ge0$.
\end{proof}

\begin{proof}[Proof of \cref{thm:teacher:main}(i)--(iv)]
Apply \cref{lem:kink:main} in the population reading, $L=R\ge L'=R_B$, at the reference point $(d+\kappa_g,H(p))$, bounded by $\pi_g$, so that $J_{\rm r}=J_g$, with $\mathcal N,\mathcal M$ and the right line of \cref{prop:teacher:kinks}. For $c>2^{-m}$ the lemma's $\Delta(c)$ is $\Lambda_2/2^m+(c-2^{-m})(4+\Lambda_1+\kappa_g)=c(4+\Lambda_1+\kappa_g)+H(p)-B_S/2^m$.

\emph{(i).} $c_{\rm hi}\ge c>0$ forces $G_1>H(p)$, hence $p<p_0<\frac14$ ($G_1$ is largest at $B=\infty$), $\Lambda_2\ge0$, and the lemma's $\Delta(c)$ is $\Delta(c)\ge0$; also $\eta_{\mathcal N}(c)\ge0$. \Cref{lem:kink:main}(i) gives $\Phi^B(c)\ge\min\{\Phi_2(c),J_g(c)-\Delta(c)\}$, with $J_g(c)-\Delta(c)=c(d-4-\Lambda_1)+B_S/2^m$, and $J_c(\pi_g)\le J_g(c)\le\Phi(c)+\Delta(c)$; $\Phi\ge\Phi^B$ as $R\ge R_B$. Finally $H(p)-B_S/2^m=(2^mH(p)-L_S+\delta_S)/2^m\le(\log(2^m+1)+\delta_S)/2^m$.

\emph{(ii).} Every $\pi\in\mathcal N$ has $J^B_c(\pi)\ge\Phi_2(c)>J_g(c)+\epsilon$, the pointwise margin of \cref{lem:kink:main}(ii); and $J_c(\pi)\le\Phi(c)+\epsilon$ implies $J^B_c(\pi)\le J_g(c)+\epsilon$, as $\Phi(c)\le J_c(\pi_g)\le J_g(c)$. So the lemma gives $\lev(\pi)=1$,
\[
\ell(\pi)\le d+\kappa_g+\frac{\Lambda_2/2^m+\epsilon}{c-2^{-m}}=d+\Delta_\ell,\qquad R_B(\pi)\ge H(p)-\frac{c\Lambda_2/2^m+2^{-m}\epsilon}{c-2^{-m}}=H(p)-\frac{c\Lambda_2+\epsilon}{c2^m-1},
\]
using $\frac\epsilon c+\frac{\epsilon/c}{c2^m-1}=\frac{\epsilon2^m}{c2^m-1}$: these are (c) and (b).
(a) By (D1), $\ell(\pi)\ge(d-4-\Lambda_1)^+$, so $R_B(\pi)=J^B_c(\pi)-c\ell(\pi)\le R_{\max}$; then $-\log(u+p+2^{-B})\le R_B$ (\cref{lem:teacher:floor}), and $\lev(\pi)=1$ gives $u>\frac34$.
(d) By (F3), $B_S\le K(\pi\mid a,K(a))+K(S\mid a,\pi)+\kappa$, and by (F4), (F3), $K(\pi\mid a,K(a))\le K(\pi)+K(a\mid\pi)-K(a)+O(1)$. By (F1), \cref{lem:teacher:list} with $t=1$, $K(a)\ge d$ and (c), this is at most $\ell(\pi)-d+4+2\log d+\kappa\le\Delta_\ell+\Lambda_3$.
(e) Given $x\in E$, $z$ is uniform on $S$ and $y^*=f_a\oplus1$, so $e(\pi)=\mathrm{avg}_SV$ with $V$ as in the proof of \cref{lem:teacher:sharpfloor}. By (d), $I(\pi)\le\Delta_\ell+\Lambda_3$, and for this budget the least admissible $j$ in that proof is $j_g$ (its threshold is $T_g$); so the $V$ side of that proof gives $\mathrm{avg}_SV\le1-u(\pi)+\gamma_g+2^{-g}$. (The argument: if $\mathrm{avg}_SV$ exceeded the least multiple $\theta$ of $2^{-g}$ above $1-u+\gamma_g$, then $S$ would lie in the set of $s_0$-subsets with average above $\theta$, which is enumerable from $(\pi,a,m,s_0,\theta)$ and, by Hoeffding's inequality for sampling without replacement \citep{hoeffding1963probability}, has at most $2^{L_S-j_g}$ elements; this would make $K(S\mid a,\pi)$ smaller than (d) allows.)

\emph{(iii).} $\lev(\pi)=1$ as in (ii). For $c2^m<1$ and $R\ge0$, $\psi_c(R)\ge cB_S+(1-c2^m)\min\{R,B_S/2^m\}$ (equality for $R\le B_S/2^m$; for $R\ge B_S/2^m$, $\psi_c(R)=R$). Compare with $J^B_c\le J_b(c)+\epsilon$ and (o), and use $L_S=B_S+\delta_S$: $(1-c2^m)\min\{R_B,B_S/2^m\}\le c(\lceil L_S\rceil+\kappa_b-B_S+4+\Lambda_1)+\epsilon=\varrho$. If $\varrho<(1-c2^m)B_S/2^m$, the minimum is $R_B\le\varrho/(1-c2^m)$. For the error label, $p\,\Ex[\min\{B,-\log\nu_\pi(y^*\mid x)\}\mid x\in E]\le R_B$, and by Jensen's inequality $\Ex[\nu_\pi(y^*\mid x)\mid E]\ge\Ex[\max\{\nu_\pi(y^*\mid x),2^{-B}\}\mid E]-2^{-B}\ge2^{-R_B/p}-2^{-B}$.

\emph{(iv).} By (D1) and $R(\pi)\ge H(p)$, $J_c(\pi)\ge c(d-4-\Lambda_1)+H(p)>c(d+\lceil L_S\rceil+\kappa_b)\ge J_c(\pi_b)$.
\end{proof}

\paragraph{Numbers.}\label{app:teacher:computations} All numbers of \cref{sec:teacher} and this appendix come from deterministic or seeded scripts in the accompanying repository. The script for \cref{tab:teacher:numbers} computes the windows in closed form and by bisection on $J_g\le\Phi_2$ (the two agree to $10^{-15}$), the identification bounds with $g\le40$, and the memoriser regime. With $\delta_S=0$, (iii) applies only for $c<0.99418\cdot2^{-m}$ at $\kappa=0$ (an illustration) and $c<0.98914\cdot2^{-m}$ at $\kappa=30$; it gives $e\ge0.954$, $0.830$, $0.658$, $0.483$, $0.222$ at $c2^m=0.5$, $0.8$, $0.9$, $0.94$, $0.97$ ($\kappa=0$) and $e\ge0.916$, $0.705$, $0.456$, $0.255$, $0.059$ ($\kappa=30$), and $e\ge\frac12$ needs $c\le0.937\cdot2^{-m}$ ($0.888\cdot2^{-m}$ at $\kappa=30$).

\subsection{Full truth tables}\label{app:teacher:fulltable}

\begin{proof}[Proof of \cref{cor:teacher:fulltable}]
On the whole truth table the empirical loss is $N$ times the population risk, so $J_\lambda(\pi;D)=N\,J_{\lambda/N}(\pi)$.
(i) Put $c=\lambda/N$, so $c2^m-1=(\lambda-2^{d-m})/2^{d-m}>0$; $J_g(c)<\Phi_2(c)$ for $c<c_{\rm hi}$, as $J_g-\Phi_2$ is strictly increasing. A $\lambda$-MAP has $J_c=\Phi(c)$, so \cref{thm:teacher:main}(ii) applies with $\epsilon=0$, and $1/(c2^m-1)=2^{d-m}/(\lambda-2^{d-m})$ gives the stated forms.
(ii) $\ell(\pi_g)+NR(\pi_g)\le d+\kappa_g+L_E+\log(N+1)$, as $NH(p)\le L_E+\log(N+1)$, and $\ell(\pi_b)+NR(\pi_b)\le d+\lceil L_E\rceil+\kappa_b$. Level-$1$ programs have $\ell+NR\ge d+L_E-\delta_E-4-\Lambda_1$ by (D2). Programs of level at least $2$ have $\ell+NR\ge NG_1$ (\cref{lem:teacher:floor}), which by the hypotheses exceeds $\ell(\pi_g)+NR(\pi_g)$, since $L_E\le NH(p)$. As $d=\log N$ and $\kappa_g$, $\kappa_b$, $\Lambda_1$ are $O(\log N)$, the claim follows.
(iii) is \cref{thm:teacher:main}(iv) at $c=\lambda/N$.
\end{proof}

At $d=m=20$, $p=2^{-10}$ and $\lambda=2$, with $g$ optimised, every $\lambda$-MAP has $u\ge0.991$, $\ell\le d+58.2$, $I\le70.8$ and $e\le0.220$ with the machine constant set to $0$ (an illustration) and $\delta_E=0$, and $\ell\le d+148.2$, $I\le190.8$, $e\le0.320$ at $\kappa=30$ (a direct computation, code in the accompanying repository); $\delta_E=20$ adds $20$ bits to $\ell$ and $I$. Near $\lambda=1$ the deficiency enters at full weight, through $\Lambda_2/(\lambda-1)$.

\subsection{Finite samples}\label{app:teacher:sampled}

\begin{lemma}[Finite sample]\status{proved}\src{steeper Lemma 2.6 (adapts T5 Lemma 1.2)}\label{lem:teacher:finite}
Let $m\le d$ and $B<\infty$. Draw $n\le2^d$ inputs uniformly without replacement and label them by $y^*$; let $\hat J_c(\pi):=c\,\ell(\pi)+\frac1n\sum_{i\le n}-\log\nu_\pi(y^*(x_i)\mid x_i)$. Let $\pi_{\rm ref}$ be a fixed program whose per-input losses lie in $[0,A]$. With probability at least $1-\delta$, every minimiser $\hat\pi$ of $\hat J_c$ satisfies
\[
J^B_c(\hat\pi)\le J_c(\pi_{\rm ref})+\epsilon_1+B\sqrt{\frac{\hat L\ln2+\ln(2/\delta)}{2n}},\qquad \epsilon_1:=A\sqrt{\frac{\ln(2/\delta)}{2n}},\quad\hat L:=\frac{J_c(\pi_{\rm ref})+\epsilon_1}{c}.
\]
\end{lemma}

\begin{proof}
Each sampled input is marginally uniform. By Hoeffding's inequality for sampling without replacement \citep{hoeffding1963probability}, with probability at least $1-\delta/2$, $\hat J_c(\pi_{\rm ref})\le J_c(\pi_{\rm ref})+\epsilon_1$. The truncated losses of each $\pi$ lie in $[0,B]$; the same inequality with confidence $(\delta/2)2^{-\ell(\pi)}$ and a union bound, which costs at most $\delta/2$ by Kraft's inequality, give, with probability at least $1-\delta/2$, simultaneously for all $\pi$,
\[
\hat J_c(\pi)\ge c\,\ell(\pi)+\frac1n\sum_{i\le n}\min\{B,-\log\nu_\pi(y^*(x_i)\mid x_i)\}\ge J^B_c(\pi)-B\sqrt{\frac{\ell(\pi)\ln2+\ln(2/\delta)}{2n}}.
\]
On both events a minimiser has $c\,\ell(\hat\pi)\le\hat J_c(\hat\pi)\le\hat J_c(\pi_{\rm ref})\le J_c(\pi_{\rm ref})+\epsilon_1$, so $\ell(\hat\pi)\le\hat L$, and the displays combine.
\end{proof}

The bound is one-sided and truncated, so no noise floor is needed; as $c$ is fixed for the given $n$, it applies to any schedule. It adapts the empirical-selection lemma of the earlier report.

\begin{proof}[Proof of \cref{cor:teacher:sampled}]
A $\lambda$-MAP with $\lambda=cn$ minimises $\hat J_c$. Apply \cref{lem:teacher:finite} with $\pi_{\rm ref}=\pi_g$, whose losses $-\log(1-p)$ and $-\log p$ lie in $[0,\log(1/p)]$: $J^B_c(\hat\pi)\le J_c(\pi_g)+\epsilon_n\le J_g(c)+\epsilon_n$.
\end{proof}

At the example's sizes, with $B=30$, $\delta=0.01$, $\delta_S=0$, $g=7$ and the machine constant set to $0$ (an illustration), a direct computation (code in the accompanying repository) gives, at $c=10^{-4}$, an excess $\epsilon_n/c$ of $92.5$ bits at $n=10^9$ ($\ell\le d+113.9$, $e\le0.284$), $9.2$ bits at $n=10^{11}$ ($u\ge0.9825$, $\ell\le d+29.8$, $e\le0.212$) and $0.9$ bits at $n=10^{13}$ ($u\ge0.9830$, $\ell\le d+21.5$, $e\le0.204$).

\subsection{The total-bits version of the example}\label{app:teacher:version1}

\begin{example}[The total-bits version of the example]\status{proved (instance of \cref{thm:teacher:main}); computed}\src{steeper Ex 2.7a; checks c6(C), c6(H)}\label{ex:teacher:version1}
Take $d=100$, $m=10$ and $p=2^{-6}$, so that $s_0=16$ and $L_S=115.6$ bits. At $n=10^4$ the model saves $n(1-H(p))=8839$ bits for its $100$ bits ($88$ bits per bit) and the error rule saves $nH(p)=1161$ bits for its $116$ bits ($10.0$ bits per bit): the sizes and rates of \cref{rem:teacher:hanni}(1). A direct computation (code in the accompanying repository) gives the window $c\in(2^{-10},c_{\rm hi}]$ with $c_{\rm hi}=5.80\cdot10^{-3}$ with the machine constant set to $0$ (an illustration) and $2.53\cdot10^{-3}$ at $\kappa=30$; at $n=10^4$ and $\kappa=0$ this is $\lambda\in(9.77,58.0]$, against $(10,100)$ for the three hypotheses. Only coarse identification is proved at this size: near-optimal hypotheses have $u>\frac34$, and at $\kappa=0$, $\delta_S=0$, $c=5.5\cdot10^{-3}$ they have $\ell\le d+28.3$ and $I\le45.6$ bits out of $115.6$; (e) is vacuous because $s_0=16$ is too small. At $n=10^4$ samples \cref{cor:teacher:sampled} gives nothing: every admissible $B$ makes its excess at least $0.321$, more than the largest margin $\Phi_2-J_g$ in the window, $0.224$.
\end{example}

\subsection{Concentration of the power posterior}\label{app:teacher:concentration}

Let $\hat J_c(\pi):=c\,\ell(\pi)+\frac1n\sum_{i\le n}-\log\nu_\pi(y_i\mid x_i)$ (\cref{rem:setup:persample}) and $\hat J_g:=\hat J_c(\pi_g)$.

\begin{proposition}[The power posterior concentrates on near-minimisers]\status{proved; computed}\src{steeper Prop 2.8; checks c5(e), c7(F); referee q1(F)}\label{prop:teacher:concentration}
Let $\lambda=cn$ with $cn\ge1$, and let $D$ be any data set of size $n$. For every $\epsilon>0$,
\[
\Pi_\lambda\big(\{\pi:\hat J_c(\pi)\ge\hat J_g+\epsilon\}\,\big|\,D\big)\big/\Pi_\lambda(\pi_g\mid D)\ \le\ 2^{-n\epsilon+(\hat J_g+\epsilon)/c}.
\]
Consequently, for data sampled as in \cref{cor:teacher:sampled}, with probability at least $1-\delta$ all but this fraction of the posterior mass lies on programs that satisfy \cref{thm:teacher:main}(ii) with $\epsilon':=\epsilon+\epsilon_1+B\sqrt{(\hat L'\ln2+\ln(2/\delta))/(2n)}$ (with $B$ as in \cref{thm:teacher:main}) in place of $\epsilon$, where $\hat L':=(J_g(c)+\epsilon_1+\epsilon)/c$, provided $J_g(c)+\epsilon'<\Phi_2(c)$.
\end{proposition}

\begin{proof}
$2^{-\lambda\ell(\pi)}\nu_\pi(D)=2^{-n\hat J_c(\pi)}$, and with the empirical risk $\hat R\ge0$, $n\hat J_c-\ell=(\lambda-1)\ell+n\hat R\ge(1-1/\lambda)\,n\hat J_c$. So on the set $\{\hat J_c\ge\hat J_g+\epsilon\}$, $2^{-n\hat J_c(\pi)}\le2^{-\ell(\pi)}2^{-(1-1/\lambda)n(\hat J_g+\epsilon)}$, and by Kraft's inequality the set has unnormalised weight at most $2^{-(1-1/\lambda)n(\hat J_g+\epsilon)}$; divide by $2^{-n\hat J_g}$ and use $n/\lambda=1/c$. For the consequence, on the event of \cref{lem:teacher:finite} with $\pi_{\rm ref}=\pi_g$, $\hat J_g\le J_g(c)+\epsilon_1$, so a program with $\hat J_c(\pi)<\hat J_g+\epsilon$ has $\ell(\pi)<\hat L'$ and $J^B_c(\pi)\le\hat J_c(\pi)+B\sqrt{(\hat L'\ln2+\ln(2/\delta))/(2n)}\le J_g(c)+\epsilon'$.
\end{proof}

Near-minimisers can be $\epsilon/c$ bits longer than the minimiser, hence the budget $\hat L'$. With $J_g=0.0232$ in place of $\hat J_g$, $c=10^{-4}$, $\epsilon=10^{-3}$, $n=10^6$ and the machine constant set to $0$ (an illustration), the ratio is at most $2^{-1000+241.7}=2^{-758.3}$. So the steeper predictive, averaged over the unseen error inputs, gives the erroneous label probability at most $\big(\sup e(\pi)+2^{-n\epsilon+(\hat J_g+\epsilon)/c}\big)\,|E|/(|E|-n)$, the supremum over the programs just described.

\subsection{\texorpdfstring{$\lambda=cn$}{lambda=cn} and weight decay}\label{app:teacher:weightdecay}

A \emph{fixed prior} is a map $w$ from programs to $[0,\infty)$ chosen before the data; its MAPs on a per-query sample $h$ maximise $w(\pi)\nu_\pi(h)$. For $\lambda\ge1$ the $\lambda$-MAPs are the MAPs of the fixed prior $2^{-\lambda\ell}$. Write $h^r$ for the sample $h$ repeated $r$ times.

\begin{proposition}[No fixed prior gives the $cn$-MAPs]\status{proved; computed}\src{imitation/decay notes-final Props 2.1--2.3, 3.2; checks d1, d3}\label{prop:teacher:nofixedprior}
Let $c>0$, and let the class be countable, with Kraft weights and kernels, as in \cref{def:split:setting}.
\begin{enumerate}[label=(\roman*),itemsep=0pt,topsep=2pt]
\item If samples $h$ and $h'$ of sizes $n$ and $n'$ have the same empirical distribution, the $cn$-MAPs of $h$ are the $cn'$-MAPs of $h'$.
\item Let $h$ have size $n_0$, let $\pi^*$ attain $P^*:=\max_\pi\nu_\pi(h)>0$, and let no $cn_0$-MAP of $h$ attain $P^*$. If $w(\pi^*)>0$, then for every $r>\max_{\hat\pi}\log\frac{w(\hat\pi)}{w(\pi^*)}\big/\log\frac{P^*}{\nu_{\hat\pi}(h)}$, the maximum over the $cn_0$-MAPs $\hat\pi$ of $h$ with $w(\hat\pi)>0$ (and every $r\ge1$ if there are none), no $crn_0$-MAP of $h^r$ is a MAP of $w$ on $h^r$. So no fixed prior with full support gives the $cn$-MAPs on all samples; nor does any fixed prior if $\pi^*$ is a $cn$-MAP of some sample whose $cn$-MAPs are not all programs. In $\Hcl{\times}$ samples as in (ii) exist for every $c>0$.
\item Per query, let $\mathcal E\ne\emptyset$ and let $w$ be summable with $w(\pi)>0$ for some $\pi\in\mathcal E$. Almost surely, for all large $n$, the MAPs of $w$ are the programs of $\mathcal E$ of largest weight. If $c>c^*$ (\cref{prop:recovery:beyond}), almost surely for all large $n$ no $cn$-MAP lies in $\mathcal E$, and every minimiser $\pi$ of $c\,\ell+R_\nu$ has $\ell(\pi)<\ell_0$ and $R_\nu(\pi)>R_\nu(\pi_0)$ for $\pi_0\in\mathcal E_0$. If $c<c^*$, the $cn$-MAPs and the MAPs of $2^{-\ell}$ are both eventually $\mathcal E_0$.
\item In the instance of \cref{def:teacher:instance}, with PQ data $\nu=\delta_{y^*}$, inputs uniform on $\{0,1\}^d$ and $B=\infty$, let $c_+:=H(p)/(B_S-4-\Lambda_1-\kappa_g)$, the denominator being positive. Every $c\in(\max\{2^{-m},c_+\},c_{\rm hi})$ has $c>c^*$.
\end{enumerate}
\end{proposition}

\begin{proof}
(i) $J_{cn}(\pi;h)/n=c\,\ell(\pi)+\Ex_{\hat P_h}[-\log\nu_\pi(y\mid x)]$ depends on $h$ only through its empirical distribution $\hat P_h$. (ii) By (i) the $crn_0$-MAPs of $h^r$ are the $cn_0$-MAPs $\hat\pi$ of $h$; they are finitely many (\cref{lem:setup:basic}), with $0<\nu_{\hat\pi}(h)<P^*$, and $w(\hat\pi)\nu_{\hat\pi}(h)^r<w(\pi^*)(P^*)^r$ for the stated $r$. A prior $w$ whose MAPs on the sample where $\pi^*$ is a $cn$-MAP are these $cn$-MAPs has $w(\pi^*)>0$: otherwise every program would have score $0$ there and be a MAP. In $\Hcl{\times}$ take $n_0$ distinct inputs whose labels $y$ have $K(y\mid\dom D)\ge n_0\log|Y|$, which exist by Kraft's inequality, with $n_0\log|Y|-\kappa-\ell_{\rm u}>\log|Y|/c$, where $\ell_{\rm u}$ is the length of a program for the uniform kernel: a lookup program attains $P^*=1$, and by \cref{lem:setup:coding} every program of likelihood $1$ has $J_{cn_0}\ge cn_0(n_0\log|Y|-\kappa)>cn_0\ell_{\rm u}+n_0\log|Y|$, the uniform kernel's value. (iii) Let $W:=\max_{\mathcal E}w$. Almost surely $\nu_\pi(h_n)=\nu(h_n)>0$ for all $\pi\in\mathcal E$ and all $n$, so a MAP $\hat\pi$ of $w$ has $w(\hat\pi)\Lambda_n(\hat\pi)\ge W$. For a cofinite set $T$ of programs, $\sum_{\pi\in T}w(\pi)\Lambda_n(\pi)$ is a nonnegative supermartingale started at $w(T)$ (\cref{lem:recovery:ratios}(a)), so by Ville's inequality some MAP lies in $T$ at some $n$ with probability at most $w(T)/W$, which tends to $0$ along cofinite sets decreasing to $\emptyset$. Of the finitely many other programs, those outside $\mathcal E$ have $\Lambda_n\to0$ almost surely (\cref{lem:recovery:ratios}(c)), and those of $\mathcal E$ with $w<W$ are never MAPs. For $c>c^*$ some $\pi'$ with $\ell(\pi')<\ell_0$ has $J_c(\pi')<J_c(\pi_0)\le J_c(\pi)$ for all $\pi\in\mathcal E$, and $J_c(\pi)>J_c(\pi_0)$ for $\pi\notin\mathcal E$ with $\ell(\pi)\ge\ell_0$; so the minimisers of $J_c$ are shorter than $\ell_0$ and lie outside $\mathcal E$, where $\KL_\mu>0$, and \cref{prop:recovery:beyond}(c) puts every $cn$-MAP among them for all large $n$. For $c<c^*$ use \cref{prop:recovery:beyond}(c) and the first part with $w=2^{-\ell}$. (iv) Here $\mathcal E$ is the set of programs with $R=0$; it contains the memoriser. Let $\pi$ minimise $J_c$, so $J_c(\pi)=\Phi(c)$. As $c2^m>1$ and $J_g(c)<\Phi_2(c)$ on $(0,c_{\rm hi})$, \cref{thm:teacher:main}(ii)(b) with $\epsilon=0$ gives $R(\pi)\ge H(p)-c\Lambda_2/(c2^m-1)$, which is positive for $c>c_+$ because $2^mH(p)-\Lambda_2=B_S-4-\Lambda_1-\kappa_g$. So no minimiser of $J_c$ lies in $\mathcal E$. For $c<c^*$ the minimisers are $\mathcal E_0$ (\cref{prop:recovery:beyond}(c)); at $c=c^*$ the programs of $\mathcal E_0$ are minimisers, since $J_c(\pi)-J_c(\pi_0)=\KL_\mu(\pi)-c^*(\ell_0-\ell(\pi))\ge0$ when $\ell(\pi)<\ell_0$. Hence $c>c^*$.
\end{proof}

Full support cannot be dropped from (ii) in general: in the class of a coin of length $1$ and a deterministic $f$ of length $1+k$, for $ck>1$ the coin is the only $cn$-MAP on every sample, and it is also the only MAP of the prior concentrated on it. With two programs that give the label $1$ probability $\frac12$ and $\frac14$ at every input, of lengths $1$ and $1+k$, the last clause of (ii) applies for every $0<ck<\log\frac32$, and a direct computation (code in the accompanying repository) finds at $ck=0.3$ no fixed prior whose MAPs are the $cn$-MAPs even on the samples of size at most $4$. At the example's sizes, with the machine constant set to $0$ or $30$ (an illustration) and $\delta_S\in\{0,20\}$, a direct computation (code in the accompanying repository) puts $c_+$ between $1.0044\cdot2^{-m}$ and $1.0115\cdot2^{-m}$.

\emph{Weight decay.} For a per-sample loss $f(\theta;z)$ let $\hat L_n:=\frac1n\sum_{i\le n}f(\cdot\,;z_i)$ and $L:=\Ex f(\cdot\,;z)$. A step $\theta\leftarrow\theta-\eta(\nabla\hat L_B(\theta)+\omega\theta)$ on a uniformly drawn minibatch $B$ is a stochastic gradient step on $F_n:=\hat L_n+\frac\omega2\|\theta\|^2$; with full batches its fixed points are the stationary points of $F_n$, for every $\eta$. If $f=-\ln p_\theta(y\mid x)$, then $nF_n(\theta)=-\ln\big(p_\theta(y_{1:n}\mid x_{1:n})\,e^{-n\omega\|\theta\|^2/2}\big)$. As in (i), the minimisers of $F_n$ depend on the sample only through its empirical distribution. If $\theta^*$ minimises $F:=L+\frac\omega2\|\theta\|^2$ and $L$ is differentiable, then $\nabla L(\theta^*)=-\omega\theta^*$, so $\theta^*$ is a stationary point of $L$ only if $\theta^*=0$; and $L(\theta^*)-L(\theta)\le\frac\omega2(\|\theta\|^2-\|\theta^*\|^2)$ for every $\theta$, the analogue of $R_\nu(\pi_c)-R_\nu(\pi)\le c(\ell(\pi)-\ell(\pi_c))$ for a minimiser $\pi_c$ of $J_c$. Unlike a vertex of the hull, a point $\theta\ne0$ is stationary for $L+\frac\omega2\|\theta\|^2$ at one value of $\omega$ at most when $L$ is differentiable; without differentiability it can be selected on an interval ($L(\theta)=|\theta-1|$ selects $\theta=1$ for every $\omega\in(0,1]$). If the decay is not multiplied by the learning rate $\alpha$ of the loss, $\theta\leftarrow\theta-\eta_t(\alpha\nabla L(\theta)+\omega_{\rm d}\theta)$, the fixed points are the stationary points of $L+\frac{\omega_{\rm d}}{2\alpha}\|\theta\|^2$: the effective coefficient then depends on $\alpha$.

\begin{example}[A worse classifier in the limit]\status{proved; computed}\src{imitation/decay notes-final Example 4.3; checks d2(A)--(D), d3(D)}\label{ex:teacher:weightdecay}
Let $y$ be uniform on $\{\pm1\}$, $\xi=1$ with probability $1-q$ and $\xi=-1$ otherwise, independently, $0<q<\frac12$, $\varepsilon>0$, $x=(\varepsilon y,\xi y)$ and $p_\theta(y\mid x)=\sigma(y\,\theta\cdot x)$. The first feature separates the classes, at scale $\varepsilon$; $\inf L=0$ is approached only by classifiers of error rate $0$. If $\varepsilon^2<1-2q$, the minimiser of $F$ has error rate $0$ for $\omega<\omega^*$ and $q$ for $\omega\ge\omega^*$ (a tie counted as an error), where
\[
\omega^*=\frac{2q\varepsilon^2}{(1-\varepsilon^2)\,a^*},\qquad a^*=\ln\frac{2(1-q)(1-\varepsilon^2)-q(1+\varepsilon^2)}{q(1+\varepsilon^2)};
\]
if $\varepsilon^2\ge1-2q$, it has error rate $0$ for every $\omega>0$. At $\varepsilon=0.1$ and $q=0.2$, $\omega^*=2.101\cdot10^{-3}$. A direct computation (code in the accompanying repository) finds the empirical minimisers with $\omega=0.01$ at error rate $0.2$ in all of 20 samples for each $n\in\{10^2,10^3,\dots,10^6\}$, and with $\omega=1/n$, the fixed prior $\mathcal N(0,I)$, at error rate $0$ in all of 20 for each $n\ge10^3$; minibatch gradient descent reaches the same regimes. If full-batch AdamW with decoupled coefficient $\omega_{\rm d}$ converges, under a non-increasing learning rate with divergent partial sums, its limit is a KKT point of $L$ under $\|\theta\|_\infty\le1/\omega_{\rm d}$ \citep{xie2024implicit}, here the minimiser of $L$ on this box, as $L$ is convex; in the same computation AdamW ends near it. That minimiser has error rate $q$ exactly when $\omega_{\rm d}\ge2\varepsilon/\ln((2-3q)/q)$, for every $\varepsilon\le1$: at $\varepsilon=0.1$, $q=0.2$ the threshold is $0.1028$, and at $\varepsilon=0.8$, $q=0.2$, where $L_2$ decay never switches, it is $0.8222$ (both to four digits).
\end{example}

\emph{Proof idea.} With $a:=\varepsilon\theta_1+\theta_2$, $b:=\varepsilon\theta_1-\theta_2$, the error rate is $(1-q)[a\le0]+q[b\le0]$. The stationarity equations of the strictly convex $F$ give $\theta_1,\theta_2>0$, so the error rate is $q[b\le0]$. At $b=0$ they force $(1-q)(1-\varepsilon^2)\sigma(-a)=\frac q2(1+\varepsilon^2)$. If $\varepsilon^2\ge1-2q$ this is impossible: for $\varepsilon\ge1$ the left side is not positive, and otherwise it needs $\sigma(-a)\ge\frac12$, against $a=2\theta_2>0$. If $\varepsilon^2<1-2q$ it fixes $a=a^*$ and $\omega=\omega^*$. As $b$ is continuous in $\omega$ and $b>0$ for $\omega\le\varepsilon^2q/\ln((2-3q)/q)$, and $\omega\,b\to(\varepsilon^2-(1-2q))/2$ as $\omega\to\infty$, the claim follows. Under the box, $\theta_1=1/\omega_{\rm d}$ because $\partial L/\partial\theta_1<0$, and for $\varepsilon\le1$ the error rate is $q$ iff $\partial L/\partial\theta_2\le0$ at $\theta_2=\varepsilon/\omega_{\rm d}$, that is iff $(1-q)\sigma(-2\varepsilon/\omega_{\rm d})\ge q/2$.

\subsection{Charging the information about the rule}\label{app:teacher:rates}

For $t\in\{2,3,\dots\}\cup\{\infty\}$ ($s(\infty):=\infty$, $\theta_{\infty-1}:=\frac12$, $G_{\infty-1}:=G_\infty$), integers $I\ge0$ and $c>0$ let
\begin{gather*}
\mathrm{Floor}_t(I):=\max\Big\{0,\ \sup_{g\ge1}\,-\log\big(p+(1-2p)\theta_{t-1}+p(2\gamma_g(I)+2^{1-g})+2^{-B}\big)\Big\},\\
\Psi_t(c):=\min_{I\in\{0,1,\dots\}}\big[cI+\mathrm{Floor}_t(I)\big],\qquad\Phi_2^\sharp(c):=\inf_tLB_t(c),\\
LB_t(c):=\max\big\{c(d-s(t)-\Lambda_1)^++G_{t-1},\ c\big(d-\min\{s(t)+2\log d,d\}-2\log(2d)-\kappa\big)+\Psi_t(c)\big\},
\end{gather*}
with $\gamma_g(I)$ as in \cref{sec:teacher:lemmas}; at $B=\infty$, $\mathrm{Floor}_t=\max\{0,F^\sharp_{t-1}\}$.

\begin{proposition}[Charging the information about the rule]\status{proved; computed}\src{split/teacher Lemma 5.1, Prop 5.2, Remark 2.6 (a2); split/verify-teacher \S6, \S8; checks t4(B), t6(B), rB7(B)}\label{prop:teacher:sharp}
In the instance of \cref{def:teacher:instance}, let $d\ge2\log(d+1)+\kappa+1$.
(i) Every program $\pi$ of level $t\in\{2,3,\dots\}\cup\{\infty\}$ has $\ell(\pi)\ge d-\min\{s(t)+2\log d,\ d\}+I(\pi)-\kappa$, and $I(\pi)\ge-2\log(2d)-\kappa$ is an integer.
(ii) Every such $\pi$ has $J^B_c(\pi)\ge LB_t(c)$ for every $c>0$. So $\Phi_2^\sharp\ge\Phi_2$ bounds $J^B_c$ on all programs of level at least $2$, and $\Phi_2^\sharp$ is continuous on $(0,\infty)$.
(iii) \Cref{thm:teacher:main}(i) and (ii), with all their conclusions, hold with $\Phi_2^\sharp$ in place of $\Phi_2$ (with $B$ as in \cref{thm:teacher:main}) for every $c$ in the window $\{c>2^{-m}:J_g(c)+\epsilon<\Phi_2^\sharp(c)\}$; there, and under the hypotheses of \cref{thm:teacher:main}(ii), (ii)(a) holds in the sharper form $u(\pi)\ge\big(2^{-R_{\max}}-p(1+2\gamma_g+2^{1-g})-2^{-B}\big)/(1-2p)$ for every integer $g\ge1$, with $\gamma_g$ as in (ii)(e).
(iv) \Cref{prop:teacher:schedules}(c) holds with $\Phi_2^\sharp$, computed at the same $B$, in place of $\Phi_2$, for every compact $[c_1,c_2]\subset\{c>2^{-m}:J_g(c)<\Phi_2^\sharp(c)\}$.
\end{proposition}

\begin{proof}[Proof of \cref{prop:teacher:sharp}]
(i) By (F3), $I(\pi)\le K(\pi\mid a,K(a))+\kappa$, and by (F4), (F3), $K(\pi\mid a,K(a))\le K(\pi)+K(a\mid\pi)-K(a)+\kappa$. Insert $K(\pi)\le\ell(\pi)+\kappa$ and $K(a\mid\pi)\le\min\{s(t)+2\log d,\ K(a)\}+\kappa$ (\cref{lem:teacher:list}, or trivially), and use $\min\{s,K(a)\}-K(a)\le\min\{s,d\}-d$, as $K(a)\ge d$. Also $K(S\mid a,\pi)\le K(S\mid a)+\kappa\le B_S+K(K(a))+\kappa$ by (F3), $K(K(a))\le2\log(K(a)+1)+\kappa$ by (F5), and $K(a)\le d+2\log(d+1)+\kappa\le2d-1$; so $I(\pi)\ge-2\log(2d)-\kappa$. $B_S$ and $K(S\mid a,\pi)$ are integers.
(ii) The first term of $LB_t$ is the bound in the proof of (o) before $s(t)\le4t$ is used, so $\Phi_2^\sharp\ge\Phi_2$ by that proof ($LB_\infty\ge G_\infty\ge\Phi_2$). For the second, with $I:=I(\pi)$, (i) gives $c\,\ell(\pi)\ge c(d-\min\{s(t)+2\log d,d\}-2\log(2d)-\kappa)+cI^+$ (if $I<0$, use $cI\ge-c(2\log(2d)+\kappa)$). \Cref{lem:teacher:sharpfloor} with the budget $I^+\ge I$, $u(\pi)\le\theta_{t-1}$ and $1-2p>0$ gives $R_B(\pi)\ge\mathrm{Floor}_t(I^+)$; add and minimise over $I^+$.
\emph{Continuity.} Let $t_c$ be the least integer with $s(t_c)\ge d-\Lambda_1$ and $s(t_c)+2\log d\ge d$. For $t\ge t_c$ the first term of $LB_t$ is $G_{t-1}$, the coefficient of $c$ in the second is $-2\log(2d)-\kappa$, and $\Psi_t$ is nondecreasing in $t$ as $\theta_{t-1}$ decreases; also $LB_\infty\ge LB_{t_c}$. So $\Phi_2^\sharp=\min_{2\le t\le t_c}LB_t$. As $p+(1-2p)\theta_{t-1}\ge\frac12$, $0\le\mathrm{Floor}_t\le1$, so on a compact subset of $(0,\infty)$ only finitely many $I$ compete in $\Psi_t$, and each $LB_t$ is continuous.
(iii) The proofs of \cref{thm:teacher:main}(i),(ii) use (o) only through ``level at least $2$ implies $J^B_c\ge\Phi_2(c)$'', that is, through $\eta_{\mathcal N}\ge\Phi_2-J_g$, and (i) uses $c\le c_{\rm hi}$ only for $J_g\le\Phi_2$; by (ii) both hold with $\Phi_2^\sharp$ on the window. Parts (a)--(e) use only level $1$, (D1), (D2) and the bounds on $\ell$ and $I$. For the sharper (a): by (ii)(d), $I(\pi)\le\Delta_\ell+\Lambda_3$, so \cref{lem:teacher:sharpfloor} with this budget, whose $j$ is $j_g$, gives $2^{-R_B}\le p+(1-2p)u+p(2\gamma_g+2^{1-g})+2^{-B}$; with $R_B\le R_{\max}$, solve for $u$. On this window $\Lambda_2\ge0$: if $\pi_g$ had level at least $2$, then by (ii) $J_g(c)\ge J_c(\pi_g)\ge J^B_c(\pi_g)\ge\Phi_2^\sharp(c)$, against $J_g(c)+\epsilon<\Phi_2^\sharp(c)$; so $u(\pi_g)=1-p>\frac34$, hence $p<\frac14$, and $\Lambda_2\ge0$ by \cref{prop:teacher:kinks}(b).
(iv) The proof of \cref{prop:teacher:schedules}(c) uses $\Phi_2$ only through \cref{thm:teacher:main}(ii) and the continuity and positivity of $\Phi_2-J_g$ on $[c_1,c_2]$.
\end{proof}

\begin{table}[!ht]
\centering\small
\begin{tabular}{@{}l|cc|cc@{}}
\toprule
$p$ & window, $\kappa=0$ & window, $\kappa=30$ & $u\ge$ & $e\le$\\
\midrule
$2^{-10}$ & $(1,\ 8.198\cdot10^{-3})$ & $(1,\ 3.686\cdot10^{-3})$ & 0.989 & 0.193\\
$1/16$ & $(1,\ 2.529\cdot10^{-3})$ & $(1,\ 6.559\cdot10^{-4})$ & 0.827 & 0.197\\
$5/64$ & $(1.014,\ 1.521\cdot10^{-3})$ & $(1.016,\ 4.285\cdot10^{-4})$ & 0.802 & 0.221\\
$3/32$ & $(2.141,\ 9.097\cdot10^{-4})$ & $(2.155,\ 2.487\cdot10^{-4})$ & 0.780 & 0.241\\
$7/64$ & $(6.223,\ 3.184\cdot10^{-4})$ & $(6.532,\ 8.558\cdot10^{-5})$ & 0.759 & 0.260\\
$15/128$ & $(21.433,\ 6.225\cdot10^{-5})$ & none & 0.750 & 0.269\\
$1/8$ & none & none & -- & --\\
\bottomrule
\end{tabular}
\caption{The windows $\{c>2^{-m}:J_g(c)<\Phi_2^\sharp(c)\}$ of \cref{prop:teacher:sharp} at $d=100$, $m=20$, $B=\infty$, $\delta_S=0$; lower ends in units of $2^{-20}$, rounded up, and upper ends rounded down. Rows with $\kappa=0$ set the machine constant to $0$ and are illustrations; $\kappa_g=2\log d+K(p)+\kappa$, with $K(p)$ taken as $2\log r$ for $p=2^{-r}$ and $2\log(h+1)+2\log(e+1)$ for $p=h/2^e$, $h>1$ odd. The last columns give the sharper (ii)(a) of \cref{prop:teacher:sharp}(iii) and the resulting (ii)(e) at $\kappa=0$, $\epsilon=0$, $g$ optimised, at $c=10^{-4}$ ($p=2^{-10}$) and $c=6.745,5.282,4.278,4.337,3.586\cdot10^{-5}$ (the other rates in order); $u$ is rounded down and $e$ up. At $p=2^{-10}$ and $c=10^{-4}$ this is the sharpest of three forms of (ii)(a): \cref{tab:teacher:numbers} uses $J_g(c)+\epsilon$ for $R_{\max}$ ($u\ge0.983$, $e\le0.200$), and \cref{rem:teacher:numbers} uses $R_{\max}$ ($u\ge0.988$, $e\le0.195$).}
\label{tab:teacher:rates}
\end{table}

\begin{proof}[Proof of \cref{prop:teacher:rates}]
Let $f(x):=-\log(\frac34-\frac p2+2p\sqrt{x\ln2/2})$, so $\psi_2(p,a)=\inf_{x\ge0}[apx+f(x)]$; at $c=a2^{-m}$ a budget $I=xs_0$ costs $cI=apx$.
\emph{The floor converges.} Let $r_g:=\delta_S+1+2\log m+2\log s_0+g+2\log g+\kappa$; in $\gamma_g(I)$, $j$ is the least integer with $j-2\log(j+1)>I+r_g$, so $I+r_g<j\le I+r_g+O(\log(I+|r_g|+2))$. Take $g=g_m\to\infty$ with $g_m=o(s_0)$; as $\delta_S=o(2^m)$, $|r_{g_m}|=o(s_0)$, so $\gamma_{g_m}(xs_0)^2=x\ln2/2+o(1)$ uniformly on compacts, and $p2^{1-g_m}\to0$. $\mathrm{Floor}_2$ is a supremum over $g$ and $\theta_1=\frac34$, so by uniform continuity $\mathrm{Floor}_2(xs_0)\ge f(x)-o(1)$ uniformly on $[0,X]$; take $X$ with $apX>\psi_2(p,a)+1$, and note $cI+\mathrm{Floor}_2\ge apX$ for $x>X$. So $\liminf_m\Psi_2(a2^{-m})\ge\psi_2(p,a)$.
\emph{All levels.} $\Psi_t\ge\Psi_2$ as $\theta_{t-1}\le\theta_1$, and the coefficient of $c$ in the second term of $LB_t$ is at least $-2\log(2d)-\kappa$, so uniformly in $t$, $J_g(c)-LB_t(c)\le c(d+\kappa_g+2\log(2d)+\kappa)+H(p)-\Psi_2(c)$. For $c\in[a2^{-m},c_m]$ the first term is at most $c_m(d+O(\log d))\to0$, and $\Psi_2(c)\ge\Psi_2(a2^{-m})$. As $a>a^*(p)$ and $\psi_2(p,\cdot)$ is nondecreasing, $\psi_2(p,a)>H(p)$, and $\epsilon:=\frac12(\psi_2(p,a)-H(p))$ works for all large $m$ (where also $d\ge2\log(d+1)+\kappa+1$).
\emph{The readings.} $\psi_2(p,\cdot)$ is an infimum of affine functions, hence concave, continuous on $(0,\infty)$ and nondecreasing. The infimum is attained at some $x^*>0$, as $f'(x)\to-\infty$ when $x\downarrow0$; so for $a'>a$, with $x^*$ the minimiser at $a'$, $\psi_2(p,a')=a'px^*+f(x^*)>apx^*+f(x^*)\ge\psi_2(p,a)$: $\psi_2(p,\cdot)$ is strictly increasing. Also $\psi_2(p,a)\le f(0)=1-\log(\frac32-p)$, and $\psi_2(p,a)\to f(0)$ as $a\to\infty$, since $f$ is continuous at $0$ and $apx\to\infty$ for $x$ bounded away from $0$. So $\psi_2(p,a)<f(0)$ for every $a$, and some $a$ has $\psi_2(p,a)>H(p)$ if and only if $f(0)>H(p)$. If $\psi_2(p,1)\ge H(p)$, strict increase gives $a^*(p)=1$; if $\psi_2(p,1)<H(p)$, continuity gives $a^*(p)>1$.
\end{proof}

\emph{Monotonicity of $h(p):=\psi_2(p,1)-H(p)$ on $(0,\frac14]$.} With $y:=\sqrt x$, $A:=\frac34-\frac p2$ and $k:=2p\sqrt{\ln2/2}$, $\psi_2(p,1)=\min_{y\ge0}[py^2-\log(A+ky)]$; the function is strictly convex with negative slope at $0$, so its minimiser $y^*$ is unique, positive, and solves $y^*(A+ky^*)=q:=1/\sqrt{2\ln2}=0.8493\ldots$. For $p\le\frac14$, $A\ge\frac58$ and $k\le0.2944$, so $y^*\le q/A\le1.359$, and $y^*\ge0.7385$ as $0.7385(\frac34+0.2944\cdot0.7385)<q$. By Danskin's theorem
\[
h'(p)=y^{*2}-\frac{2\sqrt{\ln2/2}\,y^*-\frac12}{(A+ky^*)\ln2}-\log\frac{1-p}p ,
\]
where, with $A+ky^*=q/y^*$, the middle term is $(2\sqrt{\ln2/2}\,y^*-\frac12)y^*/(q\ln2)\ge0.463$, increasing in $y^*\ge0.7385$; and $\log\frac{1-p}p\ge\log3$. So $h'(p)\le1.359^2-0.463-\log3<-0.2$. As $h(p)\ge-\log(A+ky^*)-H(p)\to\log\frac43$ for $p\to0$ and $h(\frac14)\le1-\log\frac54-H(\frac14)<0$, $h$ has exactly one root $p^*$ in $(0,\frac14)$.

At $d=100$, $m=20$, $B=\infty$ and $\delta_S=0$, a direct computation (code in the accompanying repository), with the minimum over integer budgets taken exactly, gives the windows of \cref{tab:teacher:rates}; each computed window is an interval. Of these rates only $2^{-10}$ has $G_1>H(p)$, so \cref{thm:teacher:main} gives no window for the others.
